% =====================================================================
%  TP4, Phase III, closing discussion: what the experiment does not settle,
%  and the handover to the evidential modelling of Phase IV.
%  Author: FOMETHE SOBMBANANG MAXIMILIEN
%  Fragment, included by phase3_sections.tex. It continues the section
%  opened by phase3_results.tex.
% =====================================================================

\subsection{What these experiments do not settle}
\label{sec:results-limits}

Four limits bound every claim made above, and they are stated here rather than
left for a reader to infer.

\paragraph{The images are thumbnails and the features are pixels.}
A pleural effusion is read on a radiograph as blunting of the costophrenic angle
and a fluid meniscus at the lung base, a local geometric cue occupying a small
part of the image. At \num{28}\,$\times$\,\num{28} that region covers a few
pixels, and a flat vector of \num{784} intensities gives none of the models any
notion of adjacency. The absolute scores reported here should be read as what
classical classifiers extract from a thumbnail, not as an estimate of what is
achievable on full resolution radiographs. The comparison that this work
supports is the one between families, conditions and metrics, all measured under
one fixed representation.

\paragraph{Published numbers are not directly comparable.}
The MedMNIST authors report a convolutional baseline on the same
\num{28}\,$\times$\,\num{28} ChestMNIST images~\cite{yang2023medmnist}, but
averaged over the \num{14} findings and as area under the ROC curve, whereas
this study reports a single finding and relies on average precision. A
difference in target, in metric and in the number of training images is enough
to make a direct ranking meaningless. Published figures and the figures of this
work are therefore kept in separate tables, and no claim of superiority or
inferiority is made across them.

\paragraph{The clean condition is not a clean reference.}
The annotations of ChestX-ray14, which ChestMNIST
redistributes~\cite{wang2017chestxray8}, were extracted from free-text
radiology reports by natural language processing rather than assigned by a
reader looking at the image. They therefore carry an error rate of their own,
unknown and not measurable from the data at hand. The condition called
\emph{clean} in \cref{tab:robustness} is the published annotation, not ground
truth, so the \num{5} percent corruption of Task~3.2 is an increment on top of a
baseline error whose size is not known. Any true degradation curve is shifted
along an axis whose origin is not observed, which compresses the effect that the
experiment can see.

\paragraph{One corruption model, one rate, one direction.}
The simulation flips a fixed fraction of the positive training labels to
negative, independently of the image content. Real annotation error is not
independent of the image: a faint effusion is missed more often than a large
one, which makes the noise instance dependent and concentrated near the decision
boundary, exactly where it is most harmful. The experiment also uses a single
rate and a single target finding. Results obtained under class-conditional noise
at \num{5} percent do not transfer by themselves to instance-dependent noise, to
higher rates, or to the other thirteen findings. Finally, ChestMNIST has no
explicit uncertainty category, unlike CheXpert~\cite{irvin2019chexpert}, whose
reports distinguish an uncertain mention from a negative one; the uncertainty
studied here is simulated for that reason, which is also the gap that Phase~IV
addresses with a model that represents it directly.

Two smaller points belong in the same list. Model selection uses a single
validation split rather than cross-validation, so the selection itself carries a
variance that is not quantified. And the posteriors are used only for ranking:
no calibration measure is reported, although \cref{sec:results-baselines} shows
at least one model, Gaussian naive Bayes, whose posterior is badly enough
calibrated to change which metric ranks it first.

\subsection{From constraining a model to representing evidence}
\label{sec:results-bridge}

The robustness mechanism tested in this phase is the only one classical
regularisation offers: shrink the hypothesis class so that a handful of wrong
labels cannot move the boundary far. It acts on the whole function, uniformly,
before any particular image is seen. It gives the model no way of behaving
differently on an image that resembles its training data and on one that does
not.

That limitation is visible in the outputs these models produce. A sigmoid
posterior of \num{0.5} is returned both for a radiograph that genuinely lies
between the two classes, where the evidence is balanced, and for a radiograph
unlike anything in the training set, where there is no evidence at all. The two
situations are clinically opposite and numerically identical, and no choice of
$\lambda$, of $C$ or of tree depth separates them, because a point estimate of a
probability has no room to carry that distinction. The operating point analysis
of \cref{sec:results-operating} makes the same point from another side: the
number the model emits is a score to be ranked, and the meaning attached to its
value is supplied from outside, by a threshold chosen on validation.

This is the opening that evidential deep learning fills. Rather than emitting a
probability, such a model emits the parameters of a Dirichlet distribution over
the class probabilities, so that the amount of evidence accumulated for an image
and the belief left uncommitted become explicit outputs alongside the
prediction~\cite{sensoy2018evidential}. The empirical ground established here,
that constraining capacity leaves the degradation under label noise of the same
order as the seed-to-seed variability, is what motivates looking for robustness
in the output representation rather than only in the size of the hypothesis
class. The formal treatment of that argument belongs to Phase~IV.
